\section{Pembagian Tugas}

Pembagian tugas mengikuti pembagian yang tercatat pada \texttt{README.md} repositori.

\begin{table}[H]
\centering
\begin{tabular}{@{}p{4.6cm}p{2.3cm}p{8.4cm}@{}}
\toprule
Nama & NIM & Kontribusi \\
\midrule
Muhammad Aufar Rizqi Kusuma & 13524061 & Kontrak dan fondasi bersama (model, validator, \textit{neighbor}, \textit{SearchResult}), \textit{genetic algorithm}, serta bonus multi-truk \\
Arina Azka & 13524049 & \textit{Hill climbing} wajib beserta seluruh variannya (\textit{steepest}, \textit{sideways}, \textit{stochastic}, \textit{random restart}) sebagai bonus \\
Michael James Liman & 13524106 & \textit{Simulated annealing} wajib dan \textit{objective function} alternatif (\texttt{eta\_weighted\_value} dan \texttt{compact\_value}) sebagai bonus, termasuk penjalanan \textit{hill climbing} dan \textit{genetic algorithm} pada ketiga \textit{objective function} \\
\bottomrule
\end{tabular}
\caption{Pembagian tugas anggota kelompok Three Agents}
\end{table}

\section{Repositori GitHub}

Kode sumber, data eksperimen, visualisasi, dan sumber laporan ini tersedia pada repositori berikut.

\begin{center}
\url{https://github.com/parkuskus/IF3170_G11_ThreeAgents}
\end{center}

Struktur direktori utama repositori adalah sebagai berikut.

\begin{table}[H]
\centering
\begin{tabular}{@{}p{5.2cm}p{9.8cm}@{}}
\toprule
Lokasi & Isi \\
\midrule
\texttt{src/model.py}, \texttt{validation.py}, \texttt{neighbors.py}, \texttt{objectives.py} & Representasi \textit{state}, validator \textit{constraint}, pembangkit \textit{state} awal dan \textit{neighbor}, serta ketiga \textit{objective function} \\
\texttt{src/algorithms/} & Implementasi \textit{hill climbing}, \textit{simulated annealing}, dan \textit{genetic algorithm} \\
\texttt{src/experiments/} & Pelaksana eksperimen, termasuk perbandingan \textit{objective} dan ablasi GA \\
\texttt{src/visualization/} & Pembuat plot riwayat \textit{objective} dan tampilan 3D \textit{state} \\
\texttt{examples/} & Data \textit{instance}: \texttt{ga-hard.json} (8 paket) dan \texttt{ga-demo.json} (5 paket) \\
\texttt{output/} & Hasil \textit{run} (JSON) dan visualisasi untuk setiap algoritma dan \textit{objective} (\texttt{hc-hard}, \texttt{sa-hard}, \texttt{ga-hard}, dan padanan \texttt{-demo}) \\
\texttt{tests/} & Tes fondasi, algoritma, dan eksperimen \\
\bottomrule
\end{tabular}
\caption{Struktur repositori}
\end{table}

Dependensi dipasang dengan \texttt{python -m pip install -r requirements.txt} (Python 3.10 atau lebih baru), dan seluruh tes dijalankan dengan \texttt{python -m unittest discover -s tests -v} dari direktori utama repositori. Setiap \textit{run} mencatat \textit{seed} dan parameter pada berkas \texttt{result.json} sehingga hasil pada Bab Hasil Eksperimen dan Analisis dapat direproduksi.
